La relación entre el volumen negociado y el precio de cierre evidencia comportamientos distintos entre las empresas con mayor rendimiento acumulado. TSLA presenta la mayor dispersión en el volumen de negociación, alcanzando registros superiores a los 900 millones de acciones negociadas en algunas jornadas, mientras que sus precios de cierre oscilaron hasta aproximadamente los $400. Sin embargo, la distribución de los puntos no refleja un patrón lineal claro, lo que indica que un mayor volumen negociado no implica necesariamente un mayor precio de cierre.
NVDA muestra un comportamiento más concentrado en términos de volumen, con valores que rara vez superan los 250 millones de acciones negociadas, aunque alcanza precios de cierre cercanos a los \$500. En contraste, BNTX registra volúmenes considerablemente menores, generalmente inferiores a los 20 millones de acciones, pero mantiene precios de cierre comparables a los de empresas con una actividad bursátil mucho mayor, alcanzando valores cercanos a los $400. Esto evidencia que el nivel del precio de una acción no depende directamente de la cantidad de títulos negociados durante una sesión.
En términos generales, la mayor parte de las empresas concentra sus observaciones por debajo de los 200 millones de acciones negociadas. Asimismo, FICO destaca por registrar los precios de cierre más elevados del conjunto, superando los $1,100 por acción, pese a mantener volúmenes de negociación relativamente bajos en comparación con otras compañías. En conjunto, la dispersión observada sugiere una relación débil entre el volumen negociado y el precio de cierre, indicando que ambas variables responden a factores distintos y que el incremento en el volumen de negociación no se traduce necesariamente en un incremento proporcional en el precio de la acción.